% TOP_PROBLEM: {"id": 2306071, "title": "Research Problems in Function Theory — Problem 6.71", "category_id": 8, "category": {"id": 8, "name": "analysis", "display_name": "Analysis", "description": "Limits, continuity, calculus, and function theory.", "slug": "analysis", "order_index": 8, "created_at": "2026-07-31T15:26:25.671Z"}, "status": "open"}
% FIRST_POSTED: null
% ENTRY_KIND: statement_only
% Imported from: batch03.tex; UnsolvedMath ID 2306071
\documentclass[11pt]{article}
\usepackage[margin=25mm]{geometry}
\usepackage{fontspec}
\setmainfont{Latin Modern Roman}
\usepackage{amsmath,amssymb,mathtools}
\usepackage{unicode-math}
\setmathfont{Latin Modern Math}
\usepackage{xurl,hyperref}
\hypersetup{colorlinks=true,urlcolor=blue}
\setcounter{secnumdepth}{0}
\setlength{\parindent}{0pt}
\setlength{\parskip}{5pt}
\setlength{\emergencystretch}{3em}
\newcommand{\sref}[2]{\hyperlink{source:#1:#2}{[S#2]}}
\newcommand{\cataloguescope}[1]{\paragraph{Recorded target.}#1}
\begin{document}
% UnsolvedMath ID 2306071; source catalogue code AMR-022-6071
\setcounter{section}{2306070}
\section{Research Problems in Function Theory — Problem 6.71}\label{2306071}
{\small\sffamily UnsolvedMath ID 2306071 · Analysis}
\subsection{Definitions and mathematical statement}

\paragraph{Shared notation.}$\mathbb N=\{1,2,\ldots\}$, and $\mathbb Z,\mathbb Q,\mathbb R,\mathbb C$ denote the integers, rationals, reals and complex numbers. Any other conventions are specified locally.


Write $\mathbb D=\{z\in\mathbb C:|z|<1\}$ and $\mathbb T=\partial\mathbb D$. The class $S$ consists of holomorphic injective functions $f:\mathbb D\to\mathbb C$ normalized by $f(0)=0$ and $f'(0)=1$, so $f(z)=z+\sum_{n=2}^{\infty}a_nz^n$. A function $f$ in a class $\mathcal A$ is an extreme point of that class if $f=tg+(1-t)h$, with $g,h\in\mathcal A$ and $0<t<1$, forces $g=h=f$. A normalized function $f$ is close-to-convex if there is a convex univalent function $g$ on $\mathbb D$ and a real $\eta$ such that $\operatorname{Re}(e^{i\eta}f'/g')>0$ on $\mathbb D$. A bound $I_f(R)=O((1-R)^{-1})$ means there are constants $C_f,R_f$ with $I_f(R)\le C_f/(1-R)$ for $R_f<R<1$; the constant may depend on $f$. This class-relative extreme-point convention and the distinct closed-convex-hull convention are stated separately in \sref{2306071}{3}.
\paragraph{Statement recorded in the supplied source.}
For \[I_f(R)=\int_0^{2\pi}\left|\frac{f'(Re^{it})}{f(Re^{it})}\right|^2\,dt,\] does $I_f(R)=O((1-R)^{-1})$ as $R\to1^-$ hold for every extreme point $f$ of $S$? Does the same estimate hold for every close-to-convex $f\in S$? The general recorded estimate has the additional factor $\log(1/(1-R))$. \sref{2306071}{2}
\subsection{Short English statement}
Can the extra logarithmic factor in this integral estimate be removed for extreme univalent functions and for close-to-convex functions?
\subsection{Sources}
\begin{description}
\item[\hypertarget{source:2306071:1}{S1}] ULAM.AI and the UnsolvedMath Contributors, \emph{UnsolvedMath: A Curated Collection of Open Mathematics Problems}, record 2306071 (AMR-022-6071). CC BY 4.0. \url{https://huggingface.co/datasets/ulamai/UnsolvedMath}.
\item[\hypertarget{source:2306071:2}{S2}] W. K. Hayman and E. F. Lingham, \emph{Research Problems in Function Theory} (2018), Chapter 6, Problem 6.71 and its complete update; Chapter 6 initial notation. \url{https://arxiv.org/abs/1809.07200}.
\item[\hypertarget{source:2306071:3}{S3}] V. Allu and A. Pandey, \emph{Support points of some classes of analytic and univalent functions} (2022), Introduction and Preliminaries, p.3. \url{https://arxiv.org/pdf/2209.10607}.
\end{description}
\label{end:2306071}
\end{document}
